\SourceLeaf{129}%
superficial examination, particularly to assess the practical value of the advantages we may expect from them.

The author had submitted his work to the French National Institute for assessment. The report that \textsc{Delambre} made on it and that the Institute approved is appended to the work itself. \textsc{Leonelli} was not entirely satisfied with it, and in his counterremarks, also appended, he sometimes expresses his resentment quite vigorously. We must also touch on \textsc{Delambre}'s report and \textsc{Leonelli}'s reply, as both now form part of the book itself.

The method set out in the \emph{first} part for calculating the logarithm of any given quantity $A$ rests on the idea of expressing that quantity as a product of the form
\[
10^{\mu} a \left(1+\frac{1}{10} b\right)\left(1+\frac{1}{100} c\right)\left(1+\frac{1}{1000} d\right)\ldots
\]
where $a$, $b$, $c$, $d\ldots$ denote simple integers. First, the logarithms of the individual factors will form a decreasing sequence. One can be content with a greater or smaller number of them according as the logarithm of the product is wanted to more or fewer decimal places. Second, all those individual logarithms can be kept ready in a table calculated once and for all. \textsc{Leonelli} first supplies such a table here for Briggsian logarithms to 20 decimal places. It contains first the logarithms of the integers from 2 to 9, then the logarithms for all $1+\frac{1}{10} b$, i.\,e. for 1,1; 1,2; 1,3 up to 1,9, then for all $1+\frac{1}{100} c$ and so on. Thus, besides the 8 integers, there are 9 logarithms for each order of the fractional quantities. The table need extend only to order 11$^{\mathrm{ten}}$, since the logarithms of the twelfth have the same significant digits as those of the eleventh, and so on for the succeeding orders. In this way the table includes only 107 logarithms, which suffice to calculate the logarithm of any product of the above form to 18 or 19 digits by addition alone. Digit 20$^{\mathrm{ste}}$ clearly always remains uncertain. \textsc{Leonelli} has appended a similar table of the same extent for the hyperbolic logarithms.

